% Options for packages loaded elsewhere
\PassOptionsToPackage{unicode}{hyperref}
\PassOptionsToPackage{hyphens}{url}
\documentclass[
  english,
  10pt,
  a4paper,
]{article}
\usepackage{xcolor}
\usepackage[margin=20mm,headheight=15pt]{geometry}
\usepackage{amsmath,amssymb}
\setcounter{secnumdepth}{-\maxdimen} % remove section numbering
\usepackage{iftex}
\ifPDFTeX
  \usepackage[T1]{fontenc}
  \usepackage[utf8]{inputenc}
  \usepackage{textcomp} % provide euro and other symbols
\else % if luatex or xetex
  \usepackage{unicode-math} % this also loads fontspec
  \defaultfontfeatures{Scale=MatchLowercase}
  \defaultfontfeatures[\rmfamily]{Ligatures=TeX,Scale=1}
\fi
\usepackage{lmodern}
\ifPDFTeX\else
  % xetex/luatex font selection
\fi
% Use upquote if available, for straight quotes in verbatim environments
\IfFileExists{upquote.sty}{\usepackage{upquote}}{}
\IfFileExists{microtype.sty}{% use microtype if available
  \usepackage[]{microtype}
  \UseMicrotypeSet[protrusion]{basicmath} % disable protrusion for tt fonts
}{}
\makeatletter
\@ifundefined{KOMAClassName}{% if non-KOMA class
  \IfFileExists{parskip.sty}{%
    \usepackage{parskip}
  }{% else
    \setlength{\parindent}{0pt}
    \setlength{\parskip}{6pt plus 2pt minus 1pt}}
}{% if KOMA class
  \KOMAoptions{parskip=half}}
\makeatother
\usepackage{longtable,booktabs,array}
\usepackage{caption}
\captionsetup[table]{skip=6pt}
\newcounter{none} % for unnumbered tables
\usepackage{calc} % for calculating minipage widths
% Correct order of tables after \paragraph or \subparagraph
\usepackage{etoolbox}
\makeatletter
\patchcmd\longtable{\par}{\if@noskipsec\mbox{}\fi\par}{}{}
\makeatother
% Allow footnotes in longtable head/foot
\IfFileExists{footnotehyper.sty}{\usepackage{footnotehyper}}{\usepackage{footnote}}
\makesavenoteenv{longtable}
\ifLuaTeX
\usepackage[bidi=basic,shorthands=off]{babel}
\else
\usepackage[bidi=default,shorthands=off]{babel}
\fi
\ifLuaTeX
  \usepackage{selnolig} % disable illegal ligatures
\fi
\setlength{\emergencystretch}{3em} % prevent overfull lines
\providecommand{\tightlist}{%
  \setlength{\itemsep}{0pt}\setlength{\parskip}{0pt}}
% English Muon notes; typography follows the neighboring research notes.
\usepackage{fontspec}
\usepackage{enumitem}
\usepackage{fancyhdr}
\usepackage{titlesec}
\usepackage{needspace}

\setmainfont[
  Path=/usr/local/texlive/2026/texmf-dist/fonts/opentype/public/tex-gyre/,
  UprightFont=texgyrepagella-regular.otf,
  BoldFont=texgyrepagella-bold.otf,
  ItalicFont=texgyrepagella-italic.otf,
  BoldItalicFont=texgyrepagella-bolditalic.otf
]{texgyrepagella-regular.otf}
\setsansfont[
  Path=/usr/local/texlive/2026/texmf-dist/fonts/opentype/public/tex-gyre/,
  UprightFont=texgyreheros-regular.otf,
  BoldFont=texgyreheros-bold.otf,
  ItalicFont=texgyreheros-italic.otf,
  BoldItalicFont=texgyreheros-bolditalic.otf
]{texgyreheros-regular.otf}

\definecolor{notesblue}{HTML}{174A7E}
\definecolor{notesgray}{HTML}{5D6773}
\AtBeginDocument{\hypersetup{
  colorlinks=true,
  linkcolor=notesblue,
  urlcolor=notesblue,
  pdftitle={From Norms to Optimizers: SGD and Muon}
}}

\titleformat{\section}
  {\LARGE\bfseries\sffamily\color{notesblue}}{}{0pt}{}
\titleformat{\subsection}
  {\large\bfseries\sffamily\color{notesblue}}{}{0pt}{}
\titleformat{\subsubsection}
  {\normalsize\bfseries\sffamily\color{notesgray}}{}{0pt}{}
\titlespacing*{\section}{0pt}{0pt}{1.4ex}
\titlespacing*{\subsection}{0pt}{1.8ex}{0.8ex}
\titlespacing*{\subsubsection}{0pt}{1.6ex}{0.6ex}
\setcounter{secnumdepth}{0}
\setlength{\parindent}{0pt}
\setlength{\parskip}{0.45em}
\linespread{1.06}
\setlist{itemsep=0.3em,topsep=0.4em,leftmargin=1.8em}
\emergencystretch=3em
\clubpenalty=10000
\widowpenalty=10000
\displaywidowpenalty=10000
\pagestyle{fancy}
\fancyhf{}
\fancyhead[L]{\small\sffamily\color{notesgray} Muon}
\fancyhead[R]{\small\sffamily\color{notesgray} Purshow's Notes}
\fancyfoot[C]{\small\sffamily\thepage}
\renewcommand{\headrulewidth}{0.3pt}

\newcommand{\ComparisonTableStart}{%
  \clearpage
  \begingroup
  % Use a landscape sheet so fancyhdr follows the table's orientation.
  \newgeometry{layoutwidth=297mm,layoutheight=210mm,margin=20mm,headheight=15pt}
  \setlength{\paperwidth}{297mm}
  \setlength{\paperheight}{210mm}
  \special{papersize=297mm,210mm}
  \setlength{\headwidth}{\textwidth}
  \fontsize{8.5pt}{10.5pt}\selectfont
  \setlength{\tabcolsep}{4pt}
  \renewcommand{\arraystretch}{1.12}
  \setlength{\extrarowheight}{4pt}
  \setlength{\parskip}{0pt}
  \setlength{\LTpre}{0pt}
  \setlength{\LTpost}{0pt}
}
\newcommand{\ComparisonTableEnd}{%
  \clearpage
  \endgroup
  \restoregeometry
  \special{papersize=210mm,297mm}
  \setlength{\headwidth}{\textwidth}
}
\usepackage{bookmark}
\IfFileExists{xurl.sty}{\usepackage{xurl}}{} % add URL line breaks if available
\urlstyle{same}
% fallback for those not using the hyperref driver hyperxmp:
\makeatletter
\@ifundefined{xmpquote}{\newcommand{\xmpquote}[1]{#1}}{}
\makeatother
\hypersetup{
  pdflang={en},
  hidelinks,
  pdfcreator={LaTeX via pandoc}}

\author{}
\date{}

\begin{document}

\section{From Norms to Optimizers: SGD and
Muon}\label{from-norms-to-optimizers-sgd-and-muon}

\subsection{1. A Unified View: Reducing the Loss Within an Update
Budget}\label{a-unified-view-reducing-the-loss-within-an-update-budget}

Let the current parameter matrix be \(W\in\mathbb R^{m\times n}\) and
the update be a matrix \(\Delta W\) of the same shape. The objective is
to minimize \(\mathcal L(W+\Delta W)\). Without a constraint on
\(\Delta W\), this amounts to finding the optimal parameter matrix in a
single step, which is usually difficult. Instead, we reduce the loss
gradually in the neighborhood of the current parameters: at each step,
we first set an update budget \(\varepsilon>0\), then choose an update
within that budget:

\nopagebreak[4]

\[
\min_{\|\Delta W\|\leq\varepsilon}\mathcal L(W+\Delta W).
\tag{1}
\]

The norm here measures the magnitude of the update; we leave its
specific choice open for now. When the budget is sufficiently small, a
first-order expansion of the loss around the current parameters gives:

\nopagebreak[4]

\[
\mathcal L(W+\Delta W)\approx\mathcal L(W)+\langle G,\Delta W\rangle,
\qquad G=\nabla_W\mathcal L(W).
\tag{2}
\]

Here,
\(\langle G,\Delta W\rangle=\operatorname{Tr}(G^\top\Delta W)=\sum_{i,j}G_{ij}(\Delta W)_{ij}\).
Dropping the constant \(\mathcal L(W)\), Equation (1) is approximately:

\nopagebreak[4]

\[
\min_{\|\Delta W\|\leq\varepsilon}\langle G,\Delta W\rangle.
\tag{3}
\]

To express the direction and magnitude of the update separately, write a
nonzero update as:

\nopagebreak[4]

\[
\Delta W
=\underbrace{\dfrac{\Delta W}{\|\Delta W\|}}_{\text{direction }D}
\times\underbrace{\vphantom{\dfrac{\Delta W}{\|\Delta W\|}}\|\Delta W\|}_{\text{magnitude }\rho}
=\rho D.
\tag{4}
\]

Here, \(\|D\|=1\), and \(\rho>0\) is a real scalar representing the
update magnitude, with the budget constraint \(\rho\leq\varepsilon\).
The objective is now \(\rho\langle G,D\rangle\). When \(G\ne0\), the
inner product along an optimal direction is negative, so using the full
budget along that direction yields the largest decrease in the linear
objective. We therefore only need to find an optimal unit direction and
then multiply it by the budget {[}1, 2{]}:

\nopagebreak[4]

\[
D^*\in\arg\min_{\|D\|=1}\langle G,D\rangle,
\qquad
\boxed{\Delta W^*=\varepsilon D^*.}
\tag{5}
\]

The derivations below assume a nonzero gradient; when the gradient is
zero, we may choose a zero update for the local problem.

\subsection{2. The Frobenius Norm: From Normalized Updates to
SGD}\label{the-frobenius-norm-from-normalized-updates-to-sgd}

The Frobenius norm \(\|D\|_{\mathrm F}=(\sum_{i,j}D_{ij}^2)^{1/2}\)
equals the \(\ell_2\) norm of the flattened matrix. With the Frobenius
norm, Equation (5) minimizes \(\langle G,D\rangle\) subject to
\(\|D\|_{\mathrm F}=1\). By the Cauchy-Schwarz inequality for the matrix
inner product:

\nopagebreak[4]

\[
|\langle G,D\rangle|
\leq\|G\|_{\mathrm F}\|D\|_{\mathrm F}
=\|G\|_{\mathrm F}.
\tag{6}
\]

Thus, the inner product is bounded below by \(-\|G\|_{\mathrm F}\),
attained at \(D=-G/\|G\|_{\mathrm F}\). Multiplying by the budget gives:

\nopagebreak[4]

\[
D^*=-\frac{G}{\|G\|_{\mathrm F}},
\qquad
\boxed{\Delta W^*=-\varepsilon\frac{G}{\|G\|_{\mathrm F}}.}
\tag{7}
\]

A fixed budget \(\varepsilon\) yields normalized gradient descent.
Standard SGD corresponds to a budget of
\(\varepsilon_t=\eta_t\|G_t\|_{\mathrm F}\): the learning rate
\(\eta_t\) multiplied by the Frobenius norm of the current mini-batch
gradient matrix \(G_t\). Substituting this into Equation (7) cancels the
gradient norm, giving:

\nopagebreak[4]

\[
\varepsilon_t=\eta_t\|G_t\|_{\mathrm F}
\quad\Longrightarrow\quad
\boxed{W_{t+1}=W_t-\eta_t G_t.}
\tag{8}
\]

This update can be computed elementwise:
\((\Delta W_t)_{ij}=-\eta_t G_{t,ij}\). It is also the stochastic
version of proximal gradient when there is no nonsmooth term {[}3{]}.

\subsection{3. A Fixed Spectral-Norm Budget Gives
Muon}\label{a-fixed-spectral-norm-budget-gives-muon}

Another common matrix norm is the spectral norm, which measures the
largest factor by which a matrix can amplify a vector's length (its
\(\ell_2\) norm). It equals the matrix's largest singular value:

\nopagebreak[4]

\[
\|\Delta W\|_2=\max_{\|x\|_2=1}\|\Delta W x\|_2=\sigma_{\max}(\Delta W).
\tag{9}
\]

Choosing the spectral norm and setting the budget to
\(\varepsilon=\eta\), Equation (3) minimizes
\(\langle G,\Delta W\rangle\) subject to \(\|\Delta W\|_2\leq\eta\). To
find an optimal update, first take the compact SVD of the gradient
matrix \(G\), retaining only the nonzero singular values:

\nopagebreak[4]

\[
G=U_r\Sigma_rV_r^\top
=\sum_{i=1}^r\sigma_i u_iv_i^\top,\qquad \sigma_i>0.
\tag{10}
\]

Since \(u_i,v_i\) are both unit vectors,
\(u_i^\top\Delta W v_i\geq-\|\Delta W\|_2\geq-\eta\), so:

\nopagebreak[4]

\[
\langle G,\Delta W\rangle
=\operatorname{Tr}(G^\top\Delta W)
=\sum_{i=1}^r\sigma_i u_i^\top\Delta W v_i
\geq-\eta\sum_{i=1}^r\sigma_i.
\tag{11}
\]

With \(\Delta W=-\eta U_rV_r^\top\), we have \(\|\Delta W\|_2=\eta\) and
\(u_i^\top\Delta W v_i=-\eta\) for each singular direction, exactly
attaining the lower bound. Thus, one optimal update under a fixed
spectral-norm budget is:

\nopagebreak[4]

\[
\boxed{\Delta W^*=-\eta U_rV_r^\top.}
\tag{12}
\]

This directly gives the core Muon update {[}4{]}. The original singular
values \(\sigma_i\) determine the maximum first-order decrease in the
loss, namely \(\eta\sum_i\sigma_i\); the update magnitude along each
nonzero singular direction is the same budget \(\eta\). In comparison,
standard SGD has magnitude \(\eta\sigma_i\) along these directions. In
practice, Muon approximately orthogonalizes the momentum matrix and may
also apply dimension-dependent scaling {[}5{]}.

\subsection{\texorpdfstring{4. The RMS-to-RMS Operator Norm Gives the
\(\mu\)P Version of
Muon}{4. The RMS-to-RMS Operator Norm Gives the \textbackslash muP Version of Muon}}\label{the-rms-to-rms-operator-norm-gives-the-mup-version-of-muon}

For a linear layer \(y=Wx\), \(n\) is the input dimension and \(m\) is
the output dimension. To compare the typical size of activation
components across different widths, we use the root mean square (RMS):

\nopagebreak[4]

\[
\|v\|_{\mathrm{RMS}}
=\sqrt{\frac{1}{d}\sum_{i=1}^d v_i^2}
=\frac{\|v\|_2}{\sqrt d},\qquad v\in\mathbb R^d.
\tag{13}
\]

Measuring both inputs and outputs by RMS gives the RMS-to-RMS operator
norm {[}4{]}:

\nopagebreak[4]

\[
\begin{aligned}
\|\Delta W\|_{\mathrm{RMS}\to\mathrm{RMS}}
&:=\max_{x\ne0}\frac{\|\Delta W x\|_{\mathrm{RMS}}}{\|x\|_{\mathrm{RMS}}}\\
&=\max_{x\ne0}\frac{\|\Delta W x\|_2/\sqrt m}{\|x\|_2/\sqrt n}
=\sqrt{\frac{n}{m}}\,\|\Delta W\|_2.
\end{aligned}
\tag{14}
\]

A fixed budget \(\|\Delta W\|_{\mathrm{RMS}\to\mathrm{RMS}}\leq\eta\) is
equivalent to \(\|\Delta W\|_2\leq\eta\sqrt{m/n}\). Replacing the
spectral-norm budget in Equation (12) with \(\eta\sqrt{m/n}\) therefore
yields the \(\mu\)P version of the Muon update {[}4{]}:

\nopagebreak[4]

\[
\boxed{\Delta W_{\mu\mathrm P}
=-\eta\sqrt{\frac{m}{n}}\,U_rV_r^\top.}
\tag{15}
\]

For the details of \(\mu\)P, I very, very strongly recommend these three
articles {[}6, 12, 13{]} (Su and Greg Yang really are authorities on
this).

\subsection{5. From Muon to Moonlight}\label{from-muon-to-moonlight}

Following the notation above, \(\Delta W_t\) denotes the gradient-driven
parameter increment, \textbf{including the learning rate and the minus
sign for the descent direction}; weight decay is written separately:

\nopagebreak[4]

\[
\boxed{W_{t+1}=(1-\eta_t\lambda_{\mathrm{wd}})W_t+\Delta W_t.}
\tag{17}
\]

Here, \(\lambda_{\mathrm{wd}}\) is the weight decay coefficient,
abbreviated as WD below. Let \(H_t\) be the momentum matrix passed to
orthogonalization. Standard momentum and one common formulation of
Nesterov momentum are:

\nopagebreak[4]

\[
M_t=\mu M_{t-1}+G_t,\qquad
H_t=\begin{cases}
M_t,&\text{standard momentum},\\
G_t+\mu M_t,&\text{Nesterov}.
\end{cases}
\tag{18}
\]

\nopagebreak[4]

\[
\Delta W_t=-\eta_t\begin{cases}
\operatorname{NS}_5(H_t),&\text{original Muon baseline},\\
0.2\sqrt{\max(m,n)}\,\operatorname{NS}_5(H_t),&\text{Moonlight}.
\end{cases}
\tag{19}
\]

Here, \(\operatorname{NS}_5\) denotes five steps of Newton-Schulz
approximate orthogonalization, which we will discuss in more detail
later.

\textbf{Moonlight, together with QK-Clip introduced in Kimi K2, can be
summarized in three points:}

\begin{enumerate}
\def\labelenumi{\arabic{enumi}.}
\tightlist
\item
  \textbf{Add weight decay.} Apply decoupled weight decay as in Equation
  (17) to curb the continued growth of weight and activation magnitudes.
\item
  \textbf{Align the update RMS with Adam.} Multiply by
  \(0.2\sqrt{\max(m,n)}\) so that the RMS of the update before
  multiplication by the learning rate roughly matches Adam's typical
  magnitude of \(0.2\) (as documented at \url{https://kexue.fm/}). This
  allows direct reuse of the \textbf{learning rate and weight decay
  rate} tuned for Adam, reducing the cost of another hyperparameter
  search {[}7, 8{]}.
\item
  \textbf{QK-Clip.} After each update, if the maximum attention logit of
  an attention head exceeds a threshold, scale down the Q/K weights
  proportionally to suppress attention logit explosion {[}9{]}.
\end{enumerate}

The update RMS here differs from the earlier RMS-to-RMS operator norm,
which measures the ratio of output perturbation to input magnitude. Here
it means
\(\operatorname{RMS}(\Delta W/\eta)=\|\Delta W\|_{\mathrm F}/(\eta\sqrt{mn})\),
measuring the overall size of the update's elements.

\subsection{6. Core Comparison Table}\label{core-comparison-table}

Unless otherwise noted, the table below describes pretraining
configurations. Updates follow the notation of Equation (19), and weight
decay is computed as in Equation (17).

Following Equation (19), \(W\in\mathbb R^{m\times n}\) in the table
denotes the matrix actually involved in a single update computation.
With per-head processing, \(W\) is the submatrix for one head, and
\(m,n\) are that submatrix's row and column counts. After computing
updates for each head separately, they are assembled into an update for
the full weight matrix in their original positions.

\(\operatorname{PE}_8\) denotes eight steps of Polar Express, and
\(\operatorname{SK}_{11}\) denotes 11 steps of Sinkhorn balancing with
weak-row masking; \(s\) is the update scaling factor for the current
matrix. In Muown, \(g\) is the row magnitude, \(H_{R,t}\) is the
momentum matrix for the direction variable, and
\(\widehat m_{g,t},\widehat v_{g,t}\) are Adam moment estimates of the
magnitude gradient. The MiMo formula combines the base coefficient
\(0.2\) with the additional RL scaling of \(0.5\), giving \(0.1\). See
Section 7 for the detailed computations and Equation (27) for K3's
QK-Clip.

\ComparisonTableStart

{\def\LTcaptype{none} % do not increment counter
\begin{longtable}[]{@{}
  >{\raggedright\arraybackslash}p{(\linewidth - 6\tabcolsep) * \real{0.1400}}
  >{\raggedright\arraybackslash}p{(\linewidth - 6\tabcolsep) * \real{0.3000}}
  >{\raggedright\arraybackslash}p{(\linewidth - 6\tabcolsep) * \real{0.2800}}
  >{\raggedright\arraybackslash}p{(\linewidth - 6\tabcolsep) * \real{0.2800}}@{}}
\toprule\noalign{}
\begin{minipage}[b]{\linewidth}\raggedright
Model
\end{minipage} & \begin{minipage}[b]{\linewidth}\raggedright
Core Update
\end{minipage} & \begin{minipage}[b]{\linewidth}\raggedright
Main Changes Relative to Original Muon
\end{minipage} & \begin{minipage}[b]{\linewidth}\raggedright
Parameter Allocation and Settings
\end{minipage} \\
\midrule\noalign{}
\endhead
\bottomrule\noalign{}
\endlastfoot
\textbf{Qwen3.8-Flash-Next} &
\(\Delta W_t=-0.2\eta_t\sqrt{\max(m,n)}\,\operatorname{PE}_8(H_t)\). &
\textbf{Moonlight scaling + 8-step Polar Express + finer splitting.} QKV
and GDN input projections are split by head; SwiGLU FC1 is split into
gate/up, with scaling computed from the submatrix dimensions. &
Nesterov, \(\mu=0.95\). Embedding, LM head, router, and GR low-rank
projections use AdamW; N-gram tables use Adam. \\
\textbf{DeepSeek-V4.1-Flash} &
\(\begin{aligned}[t]&\text{Backbone matrices:}\\[-1pt]&\Delta W_t=-\eta_t s\,\operatorname{NS}(H_t),\\[1pt]&\text{Engram, embedding, LM head:}\\[-1pt]&\Delta W_t=-0.18\eta_t\sqrt n\,\operatorname{SK}_{11}(H_t).\end{aligned}\)
& Adds \textbf{Head-wise Muon} on top of V4; \textbf{Engram tables,
token embedding, and LM head switch to Sinkhorn updates that balance row
and column scales}. & Muon: target RMS of 0.18, \(\mu=0.95\), WD=0.1.
Sinkhorn: \textbf{no WD}; Engram LR is additionally multiplied by 5.
Norm/non-matrix parameters use AdamW. \\
\textbf{DeepSeek-V4-Pro} &
\(\begin{aligned}[t]\Delta W_t={}&-0.18\eta_t\sqrt{\max(m,n)}\\&\quad\cdot\operatorname{HybridNS}_{8+2}(H_t).\end{aligned}\)
& \textbf{Moonlight-style scaling and WD + 10-step hybrid NS}: the first
8 steps rapidly approximate the target, and the final 2 apply a
correction; the RMS scaling coefficient is \textbf{0.18}. & Nesterov,
\(\mu=0.95\), WD=0.1. Embedding, LM head, RMSNorm, and mHC's static
bias/gating factors use AdamW. \\
\textbf{Kimi K3} & \(\Delta W_t=-\eta_t s\,\operatorname{NS}(H_t)\). &
\textbf{Per-Head Muon + QK-Clip}: Q/K/V are updated per head, followed
by clipping of Q/K weights based on each head's maximum attention logit.
& Matrix parameters use Muon, WD=0.1. \textbf{Note:} The original K3
description is rather brief; it is unclear whether it uses RMS alignment
or similar adjustments. \\
\textbf{MiMo-V2.6-Pro} &
\(\begin{aligned}[t]R_{t+1}&=R_t-0.1\eta_t\sqrt{\max(m,n)}\\&\quad\cdot\operatorname{NS}_{10}(H_{R,t}),\\[2pt]g_{t+1}&=g_t-\eta_t\frac{\widehat m_{g,t}}{\sqrt{\widehat v_{g,t}}+\epsilon},\\[2pt]\Delta W_{i,t}&=g_{i,t+1}\frac{R_{i,t+1}}{\|R_{i,t+1}\|_2}-W_{i,t}.\end{aligned}\)
& \textbf{Muown: optimize row directions and row magnitudes separately.}
Direction \(R\) uses Muon, while row magnitude \(g\) uses Adam,
explicitly controlling row norms to mitigate drift in the spectral norm
of the weights. & \textbf{Pretraining:} AdamW. \textbf{Mid-training:}
hidden-layer matrices switch to Muown; embedding, LM head, and router
retain AdamW. \textbf{RL:} LR=\(3\times10^{-6}\), WD=0, no warmup,
clip=1.0; Nesterov, \(\mu=0.95\), NS=10, with an additional factor of
0.5 on the Muon update; Adam: \(\beta_1=\beta_2=0.95\),
\(\epsilon=10^{-8}\); router frozen. \\
\textbf{Inkling} & \(W_{t+1}=(1-\eta_t\lambda_t)W_t+\Delta W_t\). &
\textbf{Weight decay is coupled to the square of the learning rate in
the schedule.} The actual fraction decayed at each step is adjusted with
\(\eta_t^2\). & \textbf{Pretraining:} large matrices use Muon; all other
parameters use Adam. \\
\textbf{MAI-Thinking-1} &
\(\Delta W_t=-\eta_t\dfrac{\widehat m_t}{\sqrt{\widehat v_t}+\delta}\),
computed elementwise. & - & Pretraining (MAI-Base-1):
\((\beta_1,\beta_2)=(0.95,0.925)\); WD is generally 0.1, with 0.01 for
attention and 0.005 for embedding. RL: \(\beta_1=\beta_2=0.95\),
\textbf{WD=0}. \\
\end{longtable}
}

\ComparisonTableEnd

\subsection{7. Detailed Formulas for the Key
Changes}\label{detailed-formulas-for-the-key-changes}

\subsubsection{7.1 Orthogonalization: NS and Its
Improvements}\label{orthogonalization-ns-and-its-improvements}

I will only give a rough outline of NS here. For the details, I very
strongly recommend Su's blog (Su really is an authority on this, lol).
Original Muon, DeepSeek-V4, and Qwen share the same computational
pattern: first normalize the momentum matrix, then iterate a matrix
polynomial. Below, \(\delta\) denotes a numerical stability constant
(set to 0 when unused), to distinguish it from the update budget
\(\varepsilon\) above:

\nopagebreak[4]

\[
X_0=\frac{H_t}{\|H_t\|_{\mathrm F}+\delta}.
\tag{20}
\]

\nopagebreak[4]

\[
\boxed{X_{k+1}=a_kX_k+b_k(X_kX_k^\top)X_k+c_k(X_kX_k^\top)^2X_k.}
\tag{21}
\]

Original Muon / Moonlight uses five iterations with fixed coefficients:

\nopagebreak[4]

\[
(a_k,b_k,c_k)=(3.4445,-4.7750,2.0315).
\tag{22}
\]

DeepSeek-V4-Pro uses Equation (22) for the first eight steps, then
switches to corrective coefficients for the last two steps to stabilize
the singular values around 1:

\nopagebreak[4]

\[
(a_k,b_k,c_k)=(2,-1.5,0.5).
\tag{23}
\]

Qwen3.8-Flash-Next uses \textbf{the step-dependent coefficient schedule
of 8-step Polar Express {[}10{]}}, with \(\delta=10^{-14}\) in the
Frobenius normalization.

\subsubsection{7.2 Splitting by Head}\label{splitting-by-head}

Qwen's splitting covers QKV, GDN input projections, and SwiGLU's
gate/up; GLM-5 applies it to MLA's Q/K/V up-projections, DeepSeek-V4.1
to Q/K, and K3 splits the Q/K/V momentum matrices.

\textbf{K3 also retains post-update weight clipping.} The general form
of K2's QK-Clip is as follows: when the maximum attention logit
\(S_{\max}^h\) of head \(h\) exceeds the threshold \(\tau\), scale down
the Q/K weights:

\nopagebreak[4]

\[
\begin{aligned}
r_h&=\begin{cases}1,&S_{\max}^h\leq\tau,\\
\tau/S_{\max}^h,&S_{\max}^h>\tau,\end{cases}\\
W_q^h&\leftarrow r_h^\alpha W_q^h,\qquad
W_k^h\leftarrow r_h^{1-\alpha}W_k^h.
\end{aligned}
\tag{27}
\]

Here, \(\alpha\) allocates the scaling between the Q and K sides. MLA
also requires distinguishing between shared and non-shared components.

\subsubsection{7.3 DeepSeek-V4.1: Balancing Row and Column Scales with
Sinkhorn}\label{deepseek-v4.1-balancing-row-and-column-scales-with-sinkhorn}

He really does love Sinkhorn\ldots{} For Engram, embedding, or the LM
head, let \(W\in\mathbb R^{m\times n}\), where \(m\) is the number of
rows (the vocabulary size for embedding and the LM head) and \(n\) is
the hidden dimension. This parameter group uses the following Nesterov
momentum formulation:

\nopagebreak[4]

\[
M_t=0.95M_{t-1}+0.05G_t,\qquad
H_t=0.95M_t+0.05G_t.
\tag{28}
\]

First, zero out rows of \(H_t\) whose norms are no greater than
\(10^{-3}\) times the average row norm, yielding \(U^{(0)}\); then
alternate between row and column normalization:

\nopagebreak[4]

\[
U^{(k)}=\begin{cases}
\operatorname{RowL2Norm}_{\delta}(U^{(k-1)}),&k\text{ is odd},\\
\operatorname{ColL2Norm}_{\delta}(U^{(k-1)}),&k\text{ is even},
\end{cases}
\quad k=1,\ldots,11.
\tag{29}
\]

At each step, divide by the corresponding row or column's \(\ell_2\)
norm plus \(\delta=10^{-20}\). The final step is row normalization.
Multiplying by \(\sqrt n\) then converts a unit row \(\ell_2\) norm into
unit row RMS, after which \(0.18\) aligns the update magnitude:

\nopagebreak[4]

\[
\boxed{W_{t+1}=W_t-0.18\eta_t\sqrt n\,U^{(11)}.}
\tag{30}
\]

This parameter group has no weight decay, and Engram's learning rate is
additionally multiplied by 5. Sinkhorn balances update magnitudes
through row and column scaling; it neither guarantees equal singular
values nor performs Muon-style orthogonalization. Its goal is to
approximately balance row and column RMS; rows that were zeroed out
remain zero.

\subsubsection{7.4 MiMo: Muown Learns Weight Row Magnitudes
Separately}\label{mimo-muown-learns-weight-row-magnitudes-separately}

Muown {[}11{]} takes the view that Muon's lack of independent control
over the magnitude of the weights themselves may cause their row norms
to keep growing. It therefore represents each row of the weight matrix
as the product of a row magnitude and a normalized direction:

\nopagebreak[4]

\[
W_i=g_i\frac{R_i}{\|R_i\|_2}.
\tag{31}
\]

Here, \(g_i\) is a learnable row magnitude. Let \(r_i=\|R_i\|_2\) and
\(d_i=R_i/r_i\). Given the weight gradient \(G_i\), the chain rule
separates the magnitude gradient \(q_i\) from the direction-variable
gradient \(F_i\):

\nopagebreak[4]

\[
q_i=\langle G_i,d_i\rangle,\qquad
F_i=\frac{g_i}{r_i}(G_i-q_id_i).
\tag{32}
\]

Use Adam to update \(g\) and Muon to update \(R\). Omitting the weight
decay term, the updates can be written as:

\nopagebreak[4]

\[
\begin{aligned}
M_t^R&=\mu M_{t-1}^R+F_t,\\
R^+&=R-\eta_ta_{m,n}\operatorname{NS}(F_t+\mu M_t^R),\\
g^+&=g-\eta_t\frac{\widehat m_g}{\sqrt{\widehat v_g}+\epsilon}.
\end{aligned}
\tag{33}
\]

Here, \(\widehat m_g,\widehat v_g\) are Adam moment estimates of \(q\);
\(g\) is updated elementwise, and \(a_{m,n}=0.2\sqrt{\max(m,n)}\). The
actual weights are then reconstructed as:

\nopagebreak[4]

\[
\boxed{W_i^+=g_i^+\frac{R_i^+}{\|R_i^+\|_2}.}
\tag{34}
\]

\textbf{References}

{[}1{]} Jianlin Su: \href{https://kexue.fm/archives/10592}{An
Appreciation of the Muon Optimizer: The Fundamental Leap from Vectors to
Matrices}.

{[}2{]} Bernstein \& Newhouse:
\href{https://arxiv.org/abs/2409.20325v2}{Old Optimizer, New Norm: An
Anthology}.

{[}3{]} Parikh \& Boyd:
\href{https://web.stanford.edu/~boyd/papers/pdf/prox_algs.pdf}{Proximal
Algorithms}.

{[}4{]} Jeremy Bernstein:
\href{https://jeremybernste.in/writing/deriving-muon}{Deriving Muon}.

{[}5{]} Keller Jordan:
\href{https://kellerjordan.github.io/posts/muon/}{Muon: An optimizer for
hidden layers in neural networks}.

{[}6{]} Yang, Simon \& Bernstein:
\href{https://papers.cool/arxiv/2310.17813}{A Spectral Condition for
Feature Learning}.

{[}7{]} Jianlin Su: \href{https://kexue.fm/archives/10739}{Muon,
Continued: Why Did We Choose to Try Muon?}.

{[}8{]} Jianlin Su: \href{https://kexue.fm/archives/11267}{Why Is Adam's
Update RMS 0.2?}.

{[}9{]} Jianlin Su: \href{https://kexue.fm/archives/11126}{QK-Clip:
Taking Muon Further Along the Road to Scaling Up}.

{[}10{]} Amsel, Persson, Musco \& Gower:
\href{https://arxiv.org/abs/2505.16932}{The Polar Express: Optimal
Matrix Sign Methods and Their Application to the Muon Algorithm}.

{[}11{]} Lion, Hübler, Li, Orvieto \& He:
\href{https://arxiv.org/abs/2605.10797v1}{Muown: Row-Norm Control for
Muon Optimization}.

{[}12{]} Jianlin Su: \href{https://kexue.fm/archives/10795}{Higher-Order
MuP: Simpler Yet More Sophisticated Spectral-Condition Scaling}.

{[}13{]} Jianlin Su: \href{https://kexue.fm/archives/11605}{Beyond MuP:
2. Linear Layers and Steepest Descent}.

\end{document}
